\section{Materiais e métodos}

\subsection{Biossensor e aquisição}
Os sensores são eletrodos impressos de carbono fabricados por serigrafia em escala piloto
\preencher{fabricante/processo, área do eletrodo de trabalho, referência e contraeletrodo}. Cada
eletrodo é medido em três etapas sucessivas da montagem do sensor, e a informação usada pelo modelo
é a mudança entre a última etapa (após o contato com a amostra) e a anterior, no mesmo eletrodo.
Dois protocolos estão previstos na pipeline: o biossensor genossensor para rRNA 16S, com carbono
limpo, sonda de captura e sonda após a incubação com a amostra, medido por voltametria de pulso
diferencial (DPV); e o imunossensor para NS1, com bloqueio, imobilização do anticorpo monoclonal
(mAb-NS1) e etapa final, medido por voltametria de onda quadrada (SWV). \preencher{protocolo
usado nos resultados, sequência e concentração da sonda ou do anticorpo}
As medições foram feitas em \preencher{eletrólito} com
ferri/ferrocianeto \preencher{concentração} em um potenciostato Metrohm Autolab controlado pelo
software NOVA 2, com faixa de potencial de \preencher{início} a \preencher{fim} V, passo de
\preencher{passo} V, amplitude de modulação de \preencher{amplitude} V e tempo de modulação de
\preencher{tempo} s.

\subsection{Amostragem e referência microbiológica}
As superfícies foram amostradas segundo a ISO~18593 \citep{iso18593} \preencher{método: swab ou
esponja, área amostrada, número de pontos, empresas e tipos de superfície}. A carga bacteriana de
referência foi obtida por plaqueamento em triplicata \preencher{meio, diluições, incubação}, e a
amostra foi rotulada como contaminada quando \preencher{critério}. A carga em
log$_{10}$\,UFC/mL foi usada como alvo da regressão.

\subsection{Leitura dos arquivos e banco de dados}
Os arquivos \texttt{.nox} gerados pelo NOVA são serializações binárias do .NET; a pipeline os lê
diretamente, sem depender do software do fabricante, implementando o formato MS-NRBF
\citep{msnrbf}. De cada medição são extraídos o potencial aplicado, a corrente diferencial
($\delta i$) e, para o controle de qualidade, o potencial efetivamente medido no eletrodo e a
corrente de base; exportações do NOVA em Excel ou CSV também são aceitas. A amostra, o eletrodo e a
etapa do sensor são deduzidos da organização das pastas, do nome do arquivo e do cabeçalho do
\texttt{.nox}, com conferência por planilha. Os dados são gravados em um banco relacional (modelo
ORM com migrações versionadas) com tabelas de empresas, coletas, amostras, plaqueamentos,
medições, curvas, controle de qualidade, modelos e previsões; cada arquivo importado é registrado
pelo seu hash SHA-256, o que garante rastreabilidade e evita duplicidade. Os colaboradores enviam
as coletas por uma plataforma web com autenticação, que grava no mesmo banco.

\subsection{Controle de qualidade}
Cada medição é reprovada automaticamente quando: tem menos de 50 pontos ou mais de 2\% de pontos
inválidos; a corrente bruta ultrapassa 5~mA (sobrecarga do canal de corrente); ou o potencial
medido se afasta do aplicado em mais de 20~mV em mais de 5\% dos pontos, descontado o
deslocamento constante da amplitude do pulso na SWV (o potenciostato não manteve o potencial). Medições com corrente máxima acima de dez vezes a mediana das medições da mesma técnica
e estágio são marcadas para revisão. O resultado é combinado com a marcação manual do laboratório,
prevalecendo a mais restritiva. Só medições aprovadas entram no treino e na previsão.

\subsection{Pré-processamento e características}
Cada voltamograma é interpolado em uma grade comum de potencial (definida no treino e reutilizada
na previsão) e suavizado por Savitzky--Golay \citep{savitzky1964} (janela 11, ordem 3).
Opcionalmente, remove-se uma linha de base linear ajustada aos 10\% de pontos de cada extremo.
Foram comparadas cinco representações (Tabela~\ref{tab:variantes}). A hipótese central é que a
\emph{diferença} entre o último e o penúltimo estágio do mesmo eletrodo isola o efeito da
hibridização e cancela a variação entre eletrodos.

\begin{table}[H]
\centering\small
\caption{Representações do sinal comparadas.}\label{tab:variantes}
\begin{tabular}{lp{10cm}}
\toprule
Variante & Descrição \\ \midrule
\texttt{dif\_resumo} & descritores de cada estágio (corrente e potencial de pico, área, largura a
meia altura) e suas diferenças entre estágios \\
\texttt{dif\_curva} & idem, mais a curva de diferença completa (último $-$ penúltimo estágio) \\
\texttt{dif\_curva\_linha\_base} & \texttt{dif\_curva} com remoção de linha de base linear \\
\texttt{dif\_curva\_sem\_suavizar} & \texttt{dif\_curva} sem suavização \\
\texttt{estagio\_final} & apenas o estágio final, sem informação dos estágios anteriores \\
\bottomrule
\end{tabular}
\end{table}

\subsection{Modelos e protocolo de validação}
Foram avaliados regressão logística (regularização L2 com três intensidades e L1), análise
discriminante linear com encolhimento, SVM com núcleo RBF (calibrada), $k$ vizinhos mais
próximos, floresta aleatória \citep{breiman2001}, \textit{extra trees}, \textit{gradient boosting}
baseado em histogramas e uma versão simplificada do SISSO \citep{ouyang2018}, todos implementados
com o scikit-learn \citep{pedregosa2011}. No SISSO simplificado, os descritores primários
padronizados geram um espaço de candidatos por uma aplicação dos operadores
$\{x, x^2, \sqrt{|x|}, \log|x|, x_i - x_j, x_i x_j, x_i/x_j\}$; a triagem por independência
segura \citep{fan2008} mantém os 25 candidatos de maior separação entre as classes (AUC) e uma
busca exaustiva escolhe a combinação de até dois candidatos com maior acurácia balanceada em uma
regressão logística.

Antes de qualquer ajuste, as amostras são divididas em três conjuntos --- treino
(\PropTreino\%), validação (\PropValidacao\%) e teste (\PropTeste\%) ---, a divisão típica
sugerida por \citet[seção~7.2]{hastie2009}: o treino ajusta os modelos, a validação escolhe entre
eles e o teste estima o erro de generalização do modelo escolhido. A divisão é feita por amostra
(todas as réplicas e estágios de uma amostra ficam no mesmo conjunto), estratificada pela classe,
para que medições da mesma amostra não apareçam em conjuntos diferentes \citep{kapoor2023}. Coletas
reservadas para validação externa (por exemplo, de outra empresa) entram inteiras no teste.

No treino, cada combinação de representação e modelo é avaliada também por validação cruzada em
cinco dobras agrupadas por amostra (\textit{StratifiedGroupKFold}), que dá a estabilidade do
desempenho e o limiar de decisão --- o que maximiza o índice J de Youden \citep{youden1950} nas
previsões fora da dobra. O modelo ajustado no treino classifica a validação com esse limiar, e a
combinação escolhida é a de maior acurácia balanceada na validação (\CriterioEscolha). Só então a
combinação é re-treinada com treino e validação juntos e avaliada uma única vez no teste, que não
participa de nenhuma escolha \citep{raschka2018,vabalas2019}. Para a regressão da carga, usam-se as
amostras contaminadas, o mesmo protocolo (escolha pelo RMSE na validação) e os modelos Ridge, PLS,
SVR, floresta aleatória, \textit{extra trees} e \textit{gradient boosting}.

Como cada amostra é medida em vários eletrodos, a decisão final por amostra usa a média das
probabilidades das réplicas aprovadas no controle de qualidade. Nas amostras classificadas como
contaminadas, a carga estimada pode ser convertida em níveis de contaminação por limites de
log$_{10}$\,UFC/mL definidos pela equipe. \preencher{limites dos níveis (baixa, moderada, alta), se
adotados}

As métricas de classificação são a acurácia balanceada (AB, média de sensibilidade e
especificidade, adequada a classes desbalanceadas), sensibilidade, especificidade e área sob a
curva ROC (AUC); para a regressão, RMSE, MAE e $R^2$ em log$_{10}$\,UFC/mL.
